\documentclass{article}
\usepackage[T1]{fontenc}
\usepackage{lmodern}
\usepackage{graphicx}
\usepackage{amsmath,amssymb}
\usepackage[a4paper,margin=2cm]{geometry}
\usepackage[absolute,overlay]{textpos}
\setlength{\TPHorizModule}{1mm}
\setlength{\TPVertModule}{1mm}
\usepackage[indonesian]{babel}
\usepackage{hyperref}
\setlength{\emergencystretch}{2em}
% unit: Halaman 24

\begin{document}

% === Halaman 24 (llm) ===

\begin{center}
\textbf{A POLYNOMIAL TIME ALGORITHM FOR BREAKING\\
THE BASIC MERKLE-HELLMAN CRYPTOSYSTEM}

by Adi Shamir

Applied Mathematics\
The Weizmann Institute\
Rehovot, Israel
\end{center}

\subsection*{Abstract}

The cryptographic security of the Merkle-Hellman cryptosystem has been a major open problem since 1976. In this paper we show that the basic variant of this cryptosystem, in which the elements of the public key are modular multiples of a superincreasing sequence, is breakable in polynomial time.

\subsection*{I. Introduction}

In 1976, Whitfield Diffie and Martin Hellman published their pioneering paper on public-key cryptography (Diffie and Hellman [1976]). The paper speculated that such cryptosystems exist and surveyed their potential applications, but did not describe actual implementations. In late 1976 and early 1977, the first public-key cryptosystems were discovered (see Merkle and Hellman [1978] and Rivest, Shamir and Adleman [1978]). Since then, many variants and a few new public-key cryptosystems have been proposed, but for a variety of reasons not claiming to give any ultimate answers in this sense.

An overview of the Merkle-Hellman cryptosystem can be found in Section II. In Section III we describe the cryptanalytic attack in an informal way, and in Section IV we analyze its performance. A discussion of the results appears in Section V.

\subsection*{II. The Basic Merkle-Hellman Cryptosystem}

The public encryption key in any Merkle-Hellman cryptosystem is a sequence of $n$ natural numbers $a_1, \dots, a_n$ (a typical value of $n$ is 100 and a typical size of each $a_i$ is 200 bits). To encrypt an $m$-bit cleartext $X = x_1 \dots x_m$ ($x_i \in \{0,1\}$), the sender computes a message-dependent partial sum of the $a_i$ elements:

\[ b = \sum_{i=1}^{n} x_i a_i \]

\end{document}
